\documentclass[10pt,a4paper]{article}
\usepackage[utf8]{inputenc}
\usepackage[T1]{fontenc}
\usepackage[spanish,es-noshorthands]{babel}
\usepackage[margin=3cm]{geometry}
\usepackage{amsmath,amssymb}
\usepackage{booktabs,tabularx,array}
\usepackage{enumitem}
\usepackage[most]{tcolorbox}
\usepackage{xcolor}
\usepackage{fancyhdr}
\usepackage{pdflscape}
\usepackage{graphicx}
\usepackage{caption}
\usepackage{subcaption}
\usepackage{float}
\usepackage[colorlinks=true,linkcolor=blue!50!black,urlcolor=blue!50!black]{hyperref}

% Permite cortar las líneas después de los guiones bajos (nombres de funciones y ficheros largos)
\let\guionbajo\_
\renewcommand{\_}{\guionbajo\allowbreak}
\emergencystretch=2em

\begin{document}
\sloppy

\begin{titlepage}
    \centering
    \vfill
    {\huge \textbf{VISIÓN ARTIFICIAL POR APRENDIZAJE PROFUNDO} \par}
    \vspace{2cm}
    {\Large Proyecto: Clasificación de lesiones de piel. Del dermatoscopio a la foto de móvil. \par}
    \vspace{1cm}
    {\Large Universidad Autónoma de Madrid \par}
    \vspace{1cm}
    {\large Iván Domínguez Hernández |  ivan.dominguezh@estudiante.uam.es \par}
    \vspace{.5cm}
    {\large Lucas Miranda López | lucas.mirandal@estudiante.uam.es \par}
    \vspace{1cm}
\end{titlepage}

\tableofcontents

\newpage


\section*{Introducción}

\noindent Hablar acerca de la introducción al problema, más adelante. 

\section{Baseline}

\noindent En esta sección se encuentra todo aquello referente a la línea base. Cosas de las que hablar:

\begin{enumerate}
    \item Separación de los datos en train, val, test: en test se deberían tener en cuenta solo datos nuevos. 
    La existencia de una fuga de datos, por tanto, supondría un problema: un mismo paciente no debería caer en 
    dos conjuntos distintos, pues de ser así no quedaría demostrado el hecho de que el modelo entrenado en train 
    generalice correctamente. 
    \begin{enumerate}
        \item Problema: HAM10000 (DermaMNIST) trae la partición hecha, y al descargar, guarda cada imagen en el directorio 
        correspondiente. DermaMNIST no guarda el identificador de lesión ni de paciente, así que no se puede comprobar si una 
        misma lesión ha caído en dos particiones. No lo podemos arreglar, solo documentarlo.
        \item El dataset de Hugging Face no trae las particiones hechas, pero sí que trae el id de cada paciente, por lo que 
        si bien las particiones tendremos que hacerlas a mano, podremos evitar la fuga de datos. 
        \item Aunque en la línea base no sea necesario, se han obtenido conjuntos de train, val y test tanto como para el dataset
        de origen como para el de destino. 
    \end{enumerate} 
    \item Etiquetas: hay más etiquetas en el conjunto de origen (7) que en el de destino (5). Todas las etiquetas de destino están 
    correspondidas con alguna etiqueta en origen, por lo que la solución que se ha pensado para esto es no tener en cuenta en origen
    aquellas que sí lo estén en destino. La intención del problema es clasificar las de destino, y de esta forma podemos hacerlo. 
    \begin{enumerate}
        \item No somos expertos en el dominio. Probablemente, lo mejor para solucionar esto sería consultar a un profesional
        y determinar cual es la mejor idea para tratar esto. Tal vez se podría encapsular según qué etiquetas dentro de otras,
        al ser un tipo concreto de lesión cutánea un tipo particular de una más global. No lo sabemos y queda pendiente en caso 
        de ampliar el proyecto en un futuro. 
        \item La codificación de las etiquetas en origen se hace con números naturales: $C \equiv \{c_i\}_{i=0}^6$. Y se hace que 
        correspondan, esas mismas etiquetas, en destino. Esta codificación es compatible con métricas que intuimos se usarán más
        adelante, como una cross-entropy, que las tendrá en cuenta como índices.
          
    \end{enumerate} 
\end{enumerate}

\end{document}